% Verbatim copy of /Users/alingalatan/workspaces/grothe/handoff/manuscript/legacy/ch25-parts/20-stability-and-smoothing.tex (file dated 2026-10-25). Only this line was added.
\subsection{Gaussian localization for a weakened pinned conjecture}
Write the odd sectors as $(T-1)\alpha+Z^2\overline\beta/\sqrt2$, with
$A=(\alpha,\beta)$ a centered jointly proper complex Gaussian vector in the
remote variables. In the convention where the B coefficient itself is
antiholomorphic, $A=(r_A,\overline{r_B})$, not $(r_A,r_B)$.
No independence between $A$ and the remote cubic $q$ is assumed.
Put $U=\|A\|^2$ and $w=\mathbb EU$.

\paragraph{Gaussian event-cost lemma (proved).}
For every event $B$ of probability $p$,
\[
 \mathbb E[U\mathbf1_B]\le wp(1-\log p),
 \qquad \mathbb E(U-d)_+\le we^{-d/w}\quad(d\ge0).
\]
The endpoint values are interpreted continuously, and $w=0$ is trivial.
Indeed $U$ has the law $\sum_j\lambda_jE_j$, with independent unit
exponentials $E_j$ and $\sum_j\lambda_j=w$. Jensen's inequality gives
$U\le_{\rm cx}wE_1$. Therefore
$\mathbb E[U\mathbf1_B]\le dp+we^{-d/w}$; take $d=-w\log p$.
For $w\le w_0\le R^2$, Markov's inequality applied to $(U-d)_+$ at
$d=R^2-w$ also gives
\[
 \mathbb P(\|A\|>R)\le e^{1-R^2/w_0},\qquad
 \mathbb E[U\mathbf1_{\{\|A\|>R\}}]
 \le w\frac{R^2}{w_0}e^{1-R^2/w_0}.
\]

\paragraph{Matrix localization implication (proved).}
Suppose a fibre expression has a global quadratic cost
$\gamma\|A\|^2$ and a sharper cost $A^*CA$ on
$\{|q|\le s_0,\ \|A\|\le R\}$, with the same baseline and linear
profile term. These majorants are hypotheses of the implication.
Let $D_+=(\gamma I-C)_+$, assume
$\mathbb P(|q|>s_0)\le p_q$ and $w\le w_0\le R^2$, and put
\[
 p_* =\min\{1,p_q+e^{1-R^2/w_0}\},\qquad
 \theta_*=p_*(1-\log p_*).
\]
Then its expected quadratic cost is at most
\[
 \mathbb E\bigl[A^*(C+\theta_*D_+)A\bigr].
\]
To prove this, the extra cost on the bad event is bounded by
$A^*D_+A$. This is another nonnegative Gaussian quadratic form, so the
event-cost lemma applies with mean $\mathbb E[A^*D_+A]$. A union bound
controls the probability of the bad event. This preserves the anisotropy
of $C$ and is valid for arbitrary dependence between $q$ and $A$.

A near-maximizer supplies an amplitude cutoff as follows. If a valid
global fibre majorant has coercivity $\delta>0$, its first-chaos content
norm is at most $c$, and the remote-only objective is at most $S+\eta_0$,
then a candidate with full objective at least $S-r$ satisfies
\[
 \sqrt w\le\frac{c+\sqrt{c^2+\delta(\eta_0+r)}}{\delta}.
\]
This follows from $\delta w-2c\sqrt w\le\eta_0+r$. The bounds $c$ and
$\delta$ must be proved uniformly in the relevant parameter cell.

\subsection{Why the local pinned matrix is not a global majorant}
At the peak profile $a=699/1250$, $b=179/1250$, write
$L=Z(a+bT)$, and define
\[
 J(q;k,\alpha,\beta)=2\mathbb E_Z
 |L+q+k+(T-1)\alpha+Z^2\overline\beta/\sqrt2|.
\]
At $q=0$, its quadratic term is exactly
\[
 \mathbb E\frac{|k+(T-1)\alpha-T\beta/\sqrt2|^2}{2|L|}.
\]
This follows by expanding the norm to second order and averaging the
visible angle. The scalar Young remainder is dominated by a multiple
of $|L|^{-1}$ times a polynomial, which is integrable, and justifies the
expansion under the integral.

For the tied multiplier $k=\alpha/3+\beta/6$, including its cost $|k|^2$,
the local matrix is
\[
 M=\mathbb E\frac{v(T)v(T)^*}{2|L|}+dd^*,\qquad
 v(T)=\binom{T-2/3}{1/6-T/\sqrt2},\quad
 d=\binom{1/3}{1/6}.
\]
Outward-rounded evaluation of its closed forms gives
\[
 M\approx\begin{pmatrix}
 .7154593425565982&-.2627445185308260\\
 -.2627445185308260&.3087284185988625
 \end{pmatrix}.
\]
The normalized curvature in direction $(2,-1)$ exceeds $.843$, so even
an infinitesimal isotropic cost $.8I$ is impossible.

More decisively, $M+I/10$ is not a global majorant. With
$q=3/4$, $\alpha=7/20$, $\beta=1/2$, and $k=1/5$, the quantity
\[
 \begin{split}
 &J(q;k,\alpha,\beta)-J(q;0,0,0)
 -2\Re\{\overline{F_0(q)}k+
          \overline{F_A(q)}\alpha+\overline{F_B(q)}\beta\}\\
 &\hspace{20mm}+|k|^2-(\alpha,\beta)^*(M+I/10)(\alpha,\beta)
 \end{split}
\]
is greater than $97/1000$. Here $F_0(q)=\mathbb E\operatorname{phase}(L+q)$
and $F_A,F_B$ are the established sector profiles. The standalone
checker \texttt{check\_false\_matrix\_majorant.py} proves this on
$256\times512$ polar cells by interval bounds. The omitted unbounded
tail is nonnegative by convexity of the norm. This refutes the proposed
majorant, not Conjecture K.

\paragraph{An unproved replacement target.}
Numerical searches support, but do not prove, a matrix remainder of the
form
\[
 J-J_0-\text{linear term}+|k|^2
 \le A^*MA+C_q|q|^2\|A\|^2+C_r\|A\|^4
\]
with $C_q=C_r=2$ for the tied multiplier. For an actual cubic, the exact
creation-annihilation identity and the existing influence bound imply
\[
 \mathbb E[|q|^2\|A\|^2]
 \le\left(v+\frac{\sqrt{3\kappa}}2\right)w,
 \qquad \mathbb E\|A\|^4\le2w^2.
\]
Thus a proof of this finite-dimensional majorant, combined with a
proved near-maximizer cutoff on $w$, would yield a useful matrix
coercivity. The majorant remains a conjectural target; neither the
searches nor the local Hessian calculation certify it.


% Self-contained section fragment. No project files are modified.
\subsection{A supremum dual from independent graded-noise stability}
\label{sec:noise-stable-weak-K}

This is a weaker certificate target, not a proof of Conjecture K or a new
numerical endpoint. Fix the visible polynomial $L$, write
\[
 \Lambda(G)=\mathbb E(|L+G|+|L-G|)-\mathbb E|G|^2/\varepsilon,
 \qquad M=\sup_{n,G\in\mathcal H_{2,1}(\mathbb C^n)}\Lambda(G).
\]
The supremum is finite by Cauchy--Schwarz and the negative quadratic price.
Let $r$ belong to one of
$\mathcal H_{1,1},\mathcal H_{2,0},\mathcal H_{1,0},\mathcal H_{0,1}$,
or the constants. An independent proper Gaussian $W$ can be assigned the
complementary grade $(1,0),(0,1),(1,1),(2,0)$, or $(2,1)$, respectively.
In the first two cases this is a genuine new Gaussian coordinate. In the
others use independent normalized sums of $X_j\overline{Y_j}$, $X_jY_j$,
or Wick cubics. These are exactly of the required grade, converge to $W$,
and have uniformly bounded second moments. Independence and uniform
integrability of the linear-growth payoff therefore prove
\[
 \Lambda(G+\sqrt{s}\,rW)\le M\qquad(s\ge0).
\]
No finite-dimensional polynomial representation of a higher-grade
Gaussian is asserted.

Put $R(a,t)=\mathbb E|a+\sqrt tW|$ and $w_\pm=L\pm G$. Define
\[
 H_{r,s}(G)=\mathbb E\sum_{\pm}
 [R(|w_\pm|,s|r|^2)-|w_\pm|]
 -\frac{s\mathbb E|r|^2}{\varepsilon}.
\]
Then exactly
\begin{equation}\label{eq:noise-near-max}
 H_{r,s}(G)\le M-\Lambda(G).
\end{equation}
Thus these are necessary cuts on near-maximizers, not universal
inequalities $H_{r,s}\le0$ for every cubic.

The heat equation with the convention $\mathbb E|W|^2=1$ gives
\[
 \partial_tR(a,t)=\tfrac14\mathbb E|a+\sqrt tW|^{-1}
 \ge\frac1{4\sqrt{a^2+t}},
\]
where Jensen is applied to $x\mapsto x^{-1/2}$. Integration proves
\[
 R(a,t)-a\ge\tfrac12(\sqrt{a^2+t}-a)
 =\frac{t}{2(\sqrt{a^2+t}+a)}.
\]
Consequently the nonsingular algebraic gain
\[
 B_{r,s}(G)=\frac12\mathbb E\sum_{\pm}
 [\sqrt{|w_\pm|^2+s|r|^2}-|w_\pm|]
 -\frac{s\mathbb E|r|^2}{\varepsilon}
\]
satisfies $B_{r,s}\le H_{r,s}\le M-\Lambda(G)$. At a
$\delta$-near-maximizer, this also yields
\[
 \mathbb E\sum_{\pm}
 \frac{|r|^2}{\sqrt{|w_\pm|^2+s|r|^2}}
 \le\frac{4\mathbb E|r|^2}{\varepsilon}+\frac{4\delta}{s}.
\]
The convention is $0/0=0$; the quotients are bounded by
$|r|/\sqrt s$, so are integrable.

\paragraph{Supremum-dual theorem.}
Choose finitely many gains $C_j$, each an established lower bound on an
admissible $H_{r_j,s_j}$. Let $\lambda_j(G)\ge0$ and suppose
$\sum_j\lambda_j(G)\le A<\infty$ uniformly. Then
\begin{equation}\label{eq:noise-sup-dual}
 M\le\sup_G\left\{\Lambda(G)-\sum_j\lambda_j(G)C_j(G)\right\}.
\end{equation}
Indeed, take $\Lambda(G_n)\to M$ and set
$\delta_n=M-\Lambda(G_n)$. Equation~\eqref{eq:noise-near-max}
gives $C_j(G_n)\le\delta_n$, so the expression inside the supremum is
at least $\Lambda(G_n)-A\delta_n\to M$. This needs neither attainment
nor differentiation under an expectation. A finite boxwise choice of
bounded nonnegative multipliers is allowed. The positive noise price
$\lambda_j s_j\mathbb E|r_j|^2/\varepsilon$ in the dual must be retained.

At an attained global maximizer in the stable closure, Fatou as $s\downarrow0$
additionally gives
\[
 \mathbb E\left[|r|^2
 (|L+G|^{-1}+|L-G|^{-1})\right]
 \le4\mathbb E|r|^2/\varepsilon.
\]
For an orthonormal finite family in any one admissible grade this is a
matrix inequality with right side $(4/\varepsilon)I$. Constant and
remote directions independent of the diffuse benchmark buffer saturate
the infinitesimal constraint at the Gaussian optimum. Visible quadratic
directions generally do not. Finite-step cuts, rather than an unproved
existence of a maximizer, are the rigorous certificate interface.

\paragraph{Parity obstruction for a different dual.}
The exact constraint $Ph=tZ$ permits free multipliers in the remote first
chaos, but these do not correct the pure shared $A/B$ baseline.
There $G(-Z,X)=G(Z,X)$ and $L(-Z)=-L(Z)$, so
\[
 h_0=\tfrac12\{\operatorname{phase}(L+G)+
                   \operatorname{phase}(L-G)\}
\]
is odd in $Z$. Hence $\mathbb E_Zh_0=0$ and all remote first-chaos
coefficients already vanish. Such multipliers may still help even-sector
interactions. Neither this observation nor the noise-dual theorem above
supplies the remaining all-$G$ majorant or a compatible flank cover.



\subsection{A fixed smoothing target and the full fresh-coordinate response operator}
\label{sec:k-fixed-smoothing-response}

\paragraph{Provenance (IDE note, 20 October).} Lower-bound GPT stability-geometry package (\code{conjK-stability-geometry-20261020.zip}, archived in
\pathref{gpt-reviews/stability-geometry-20261020/}), ported with three typographical fixes. IDE replay \code{python3 -I verify\_all.py}: all algebra audits PASS
(25+25 fresh-coordinate moments, 128 Hermite lowering identities, ascent completion, fourth-order optimisation, budget fraction, 12 hashes). IDE checks \st{P}: the
Jensen smoothing step, $pg_0^2=625/7745536$, $a_g\le\sqrt\pi/(2g)$, and the fourth-order constant $2\cdot\frac1{24}\cdot\frac{3\sqrt\pi}{4g^3}\cdot23<\frac{23}{9g^3}$.
The 23 is a fourth-moment constant for the probes $VA+\bar VB$ only, not for all $\Hc21$ cubics (the Wick cubic has $44$). The matrix
$\bigl(\begin{smallmatrix}1&-1\\-1&1\end{smallmatrix}\bigr)$ is only a linear-algebra obstruction, not a cubic example.

\paragraph{Status and scope.}
The statements below are analytic proofs [P]. The accompanying standard-library
checker audits their finite Gaussian algebra and rational constants [C]; it is
not a numerical certificate for Conjecture K. No endpoint changes. The visible
polynomial is fixed, $L=Z(a+b|Z|^2)$, and $p=1/\varepsilon>0$.
The original coordinates are $(Z,X_1,\ldots,X_n)$, independent proper standard
complex Gaussians. An arbitrary $F\in\mathcal H_{2,1}$ may involve all these
coordinates and every visible sector. $U$ and $V$ below are fresh, mutually
independent proper Gaussians, independent of the original coordinates.

For a fixed $g>0$ define
\[
 J_g(F)=\mathbb E|L+F+gU|+\mathbb E|L-F+gU|
          -p(\|F\|_2^2+g^2),\qquad
 M_g=\sup_{n,F\in\mathcal H_{2,1}}J_g(F).
\]
The same $U$ may be used in these two expectations: its symmetry shows that
$J_g(F)=\Lambda(F+gU)$. This is a definition in the Gaussian-buffer extension;
$gU$ is not asserted to be a finite-dimensional third-chaos polynomial.
$M_g$ is finite by the triangle inequality and the negative quadratic price.

\subsubsection{The zero-buffer target can be replaced by one positive buffer [P]}
Conditional Jensen gives, for every square-integrable $F$ independent of $U$,
\begin{equation}\label{eq:sg-smoothing}
 \Lambda(F)\le J_g(F)+pg^2.
\end{equation}
Consequently $M_0\le M_g+pg^2$. Conversely $M_g\le M_0$ when $M_0$
is the supremum over all finite pure cubics: replace $U$ by normalized sums
of independent remote Wick cubics. The central limit theorem, independence,
and uniform integrability of the linear-growth payoff justify the limit;
the variance price is unchanged.

At the peak $p=2500/1891$, take $g_0=1/128$. Then
\[
 pg_0^2=\frac{625}{7745536}=0.000080691639674775\ldots,
\]
\[
 \eta_{AB}-pg_0^2
 =\frac{843807243}{84837824000000}
 =0.000009946120765662\ldots.
\]
Thus the following single, smooth, dimension-free statement is sufficient
for the stipulated peak budget:
\begin{equation}\label{eq:sg-new-target}
 \boxed{M_{1/128}\le S+\frac{99}{10^7}.}
\end{equation}
\emph{IDE correction (20 October).} $\eta_{AB}$ is the room above the certified remote cap $C'=2747927/2000000\approx S+9.97\cdot10^{-6}$, not above $S$
(\autoref{ch:remote-cubic-refinement}, budget paragraph). The sufficient target is therefore the weaker
$M_{1/128}\le C'+\eta_{AB}-625/7745536\approx S+1.992\cdot10^{-5}$, about twice the allowance in \eqref{eq:sg-new-target}. The statement below remains correct as a sufficient
condition.

It would imply $M_0<S+\eta_{AB}$. Equation~\eqref{eq:sg-new-target} is
\textbf{open}. This arithmetic does not pay separately for a remote error:
the $99/10^7$ allowance is for the \emph{entire} smoothed objective.
Profile uniformity and the complete row cover remain necessary for an endpoint.
The Jensen observation is elementary; the contribution here is its explicit
budget and its use with the nonsingular operator below.

\subsubsection{Exact ascent and a dimension-free stationarity error [P]}
Write $Q$ for the orthogonal projection onto the full original
$\mathcal H_{2,1}$ and set
\[
 k_g(F)=\tfrac12\mathbb E_U\{\operatorname{phase}(L+F+gU)
                         -\operatorname{phase}(L-F+gU)\},
 \qquad r_g=Qk_g(F)-pF.
\]
All scalar products are conjugate-linear in the first argument. The unpriced
part of $J_g$ is a convex differentiable function of $F$, with real derivative
$2\operatorname{Re}\langle k_g(F),H\rangle$. Its supporting hyperplane gives
\[
 J_g(F+H)-J_g(F)\ge
       2\operatorname{Re}\langle r_g,H\rangle-p\|H\|_2^2
       \quad(H\in\mathcal H_{2,1}).
\]
Taking $H=r_g/p$ proves the global ascent bound
\begin{equation}\label{eq:sg-ascent}
 \boxed{J_g(Qk_g(F)/p)-J_g(F)\ge\|r_g\|_2^2/p.}
\end{equation}
In particular, a $\delta$-near-maximizer of $M_g$ satisfies
$\|r_g\|_2^2\le p\delta$. No attainment, small-amplitude assumption,
large buffer, or independence between the sectors is used.

\subsubsection{The missing off-diagonal noise block [P]}
Put $w_\pm=L\pm F+gU$ and, taking only the $U$ expectation, define
\[
 a_g=\tfrac14\mathbb E_U(|w_+|^{-1}+|w_-|^{-1}),\qquad
 b_g=\tfrac14\mathbb E_U
       (\overline{w_+}^{\,2}|w_+|^{-3}
        +\overline{w_-}^{\,2}|w_-|^{-3}).
\]
These are functions of the original coordinates. Gaussian convolution gives
$|b_g|\le a_g\le\sqrt\pi/(2g)$: the inverse-distance potential of a
proper Gaussian is radial decreasing and its value at zero is $\sqrt\pi/g$.
This also follows by writing both radial kernels as positive superpositions
of disks. Thus all following compressed operators are bounded.

On $\mathcal V=\mathcal H_{1,1}\oplus\mathcal H_{0,2}$ let
\begin{equation}\label{eq:sg-operator}
 \mathcal T_{F,g}=
 \begin{pmatrix}\Pi_{1,1}&0\\0&\Pi_{0,2}\end{pmatrix}
 \begin{pmatrix}a_g&-\overline{b_g}\\-b_g&a_g\end{pmatrix}
 \bigg|_{\mathcal V}.
\end{equation}
It is complex-linear, self-adjoint, positive semidefinite, and has norm
at most $\sqrt\pi/g$.
For $A\in\mathcal H_{1,1}$, $B\in\mathcal H_{2,0}$, the admissible cubic
$H=VA+\bar V B$ has, conditionally on the old coordinates,
\[
 \mathbb E_V|H|^2=|A|^2+|B|^2,
 \qquad \mathbb E_VH^2=2AB.
\]
Taylor expansion at zero therefore gives
\begin{equation}\label{eq:sg-joint-second}
 J_g(F+tH)-J_g(F)
 =t^2\langle(A,\bar B),(\mathcal T_{F,g}-pI)(A,\bar B)\rangle+O(t^4).
\end{equation}
Odd terms vanish under $V\mapsto-V$. The mixed term is
$-2\operatorname{Re}\mathbb E[b_gAB]$. Separate proper-noise tests on
$A$ and $B$ see only the diagonal blocks, so do not establish
$\mathcal T_{F,g}\preceq pI$.

For orientation, the constant Hermitian matrix
$\left(\begin{smallmatrix}1&-1\\-1&1\end{smallmatrix}\right)$ has each
diagonal smaller than the peak $p$ but largest eigenvalue $2>p$.
This is only an exact \emph{linear-algebra} obstruction to a diagonal-only
inference, not a realizable cubic example or a counterexample to K.

\paragraph{Completeness for one fresh coordinate.}
Every cubic increment orthogonal to the old-coordinate sigma-field has
the unique form
\[
 H=VA+\bar V B+(|V|^2-1)C+V^2D/\sqrt2+H(V)E,
\]
with old grades $(1,1),(2,0),(1,0),(0,1),(0,0)$, respectively, and
$H(V)=V(|V|^2-2)/\sqrt2$. Orthogonality gives conditional variance
$|A|^2+|B|^2+|C|^2+|D|^2+|E|^2$ and conditional pseudovariance
$2AB+C^2$. Hence the entire quadratic coefficient is
\[
 \mathbb E[(a_g-p)(|A|^2+|B|^2+|C|^2+|D|^2+|E|^2)
               -\operatorname{Re}\{b_g(2AB+C^2)\}].
\]
Besides the joint block~\eqref{eq:sg-operator}, there is a real-linear
$C$ block and proper scalar $D,E$ blocks. This statement completes the
fresh-coordinate quadratic test, not the classification of stable cubics.

\subsubsection{Quantitative spectral pinning of the remote derivatives [P]}
For each remote coordinate $X_j$ put
\[
 d_j=\partial_{X_j}F\in\mathcal H_{1,1},\quad
 f_j=\partial_{\bar X_j}F\in\mathcal H_{2,0},\quad
 v_j=(d_j,\overline{f_j})\in\mathcal V.
\]
Then the following is an \emph{identity for every} $F$:
\begin{equation}\label{eq:sg-pinning}
 \boxed{(\mathcal T_{F,g}-pI)v_j
       =(\partial_{X_j}r_g,\overline{\partial_{\bar X_j}r_g}).}
\end{equation}
Indeed differentiation of the smoothed phase gives
\[
 \partial_{X_j}k_g=a_gd_j-\overline{b_g}\,\overline{f_j},\qquad
 \overline{\partial_{\bar X_j}k_g}
       =a_g\overline{f_j}-b_gd_j.
\]
There are no derivatives of $L$ in these equations because $X_j$ is remote.
Hermite lowering gives
$\partial_X Q=\Pi_{1,1}\partial_X$ and
$\partial_{\bar X}Q=\Pi_{2,0}\partial_{\bar X}$.
The functions are in Gaussian Sobolev space: smoothing bounds the phase
derivative by a constant times $1/g$, while derivatives of $F,L$ are
polynomials. Hermite lowering may therefore be justified first on polynomial
approximations and then in Sobolev norm. These observations prove the identity.

The cubic derivative isometry gives
\[
 \sum_{j\ \text{remote}}\bigl(\|\partial_{X_j}r_g\|_2^2+
                   \|\partial_{\bar X_j}r_g\|_2^2\bigr)
 \le3\|r_g\|_2^2.
\]
Together with~\eqref{eq:sg-ascent}, it follows that every
$\delta$-near-maximizer of $M_g$ satisfies
\begin{equation}\label{eq:sg-quantitative-pinning}
 \boxed{\sum_j\|(\mathcal T_{F,g}-pI)v_j\|_2^2\le3p\delta.}
\end{equation}
At an exact stationary point all nonzero remote derivative vectors are
eigenvectors with eigenvalue exactly $p$. This is the infinitesimal form
of the symmetry rotating an old remote coordinate into a fresh coordinate.
It is not an independent hypothesis added to stationarity. Its value is that
it expresses that condition in the same operator as the noise-stability test.
Unlike a bound through fifth derivatives of a radial profile, the error
constant $3p$ has no inverse powers of $g$.

\paragraph{A useful exclusion consequence.}
If $\mathcal T_{F,g}\preceq(p-d)I$ for some $d>0$, then
\[
 \|F-\mathbb E[F\mid Z]\|_2^2\le 3p\delta/d^2.
\]
Each Hermite monomial involving a remote variable contributes at least
its variance to the sum of remote derivative energies, which proves the
claim from~\eqref{eq:sg-quantitative-pinning}. This conditional implication
does not claim that such a gap holds uniformly; at the diffuse benchmark
the relevant spectrum reaches $p$.

\subsubsection{Quantitative stability and a zero-buffer finite-step substitute [P]}
For a unit vector $(A,\bar B)\in\mathcal V$ set
$\lambda=\langle(A,\bar B),(\mathcal T_{F,g}-pI)(A,\bar B)\rangle$.
For these probes one can improve the generic cubic fourth-moment constant
from $44$ to $23$. Put $u=\|A\|_2^2$, $v=\|B\|_2^2$. The complex
quadratic Gaussian moment formulas give
$\mathbb E|A|^4\le9u^2$ and $\mathbb E|B|^4\le6v^2$.
Indeed, writing $A=X^*MX-\operatorname{tr}M$ gives $u=\operatorname{tr}(MM^*)$
and
\[
 \mathbb E|A|^4=2u^2+|\operatorname{tr}M^2|^2+
 4\|M^2\|_{\rm HS}^2+2\operatorname{tr}((MM^*)^2)\le9u^2.
\]
For $B=X^{\mathsf T}NX$ with $N=N^{\mathsf T}$,
$v=2\operatorname{tr}(NN^*)$ and
$\mathbb E|B|^4=2v^2+16\operatorname{tr}((NN^*)^2)\le6v^2$.
These are Wick pairings: for the first formula the pairings contribute
$2u^2+|\operatorname{tr}M^2|^2$, and the six cyclic fourth cumulants
contribute the last two terms. Conditional Gaussian integration now gives
\[
 \mathbb E|VA+\bar VB|^4
 =2\mathbb E|A|^4+8\mathbb E|AB|^2+2\mathbb E|B|^4
 \le18u^2+60uv+12v^2\le23(u+v)^2.
\]
The first inequality uses $24\sqrt6<60$. For the second, the difference is
$5u^2-14uv+11v^2=((5u-7v)^2+6v^2)/5$.
Thus no generic cubic fourth-moment input is needed for this bound.
For $\Phi_g(w)=2\mathbb E|w+gU|$, averaging one-dimensional absolute values
over planar directions yields
\[
 |D_e^4\Phi_g(w)|\le\frac{3\sqrt\pi}{2g^3}\quad(|e|=1).
\]
For completeness, the one-dimensional smoothed absolute value has fourth
derivative $4(2x^2/g^2-1)e^{-x^2/g^2}/(g^3\sqrt\pi)$ in absolute value
at most $4/(g^3\sqrt\pi)$. The angular fourth moment is $3/8$ and
$2|w|=\pi\mathbb E_\theta|\operatorname{Re}(e^{-i\theta}w)|$.
Taylor's theorem, symmetrization and $\sqrt\pi<16/9$ give
\[
 J_g(F+tH)-J_g(F)\ge\lambda t^2-\frac{23}{9g^3}t^4.
\]
Thus
\begin{equation}\label{eq:sg-spectral-stability}
 \boxed{M_g-J_g(F)\ge\frac{9g^3}{92}
           (\lambda_{\max}(\mathcal T_{F,g})-p)_+^2.}
\end{equation}
This is valid over the entire finite quadratic feature space, in every
dimension. Its constant becomes small as $g\downarrow0$: it must not be
advertised as a uniform zero-buffer stability estimate.

There is also a nonsingular finite-step lower gain without smoothing.
For $w\in\mathbb C$, $H\in\mathbb C$, and $t>0$,
\[
 \frac{|w+tH|+|w-tH|}{2}-|w|
 \ge\frac{t^2(\operatorname{Im}(\bar wH))^2}
             {2(|w|+t|H|)^3},
\]
with value zero on the right at $(w,H)=(0,0)$. Integrate the second
derivative along the line segment; away from its possible zero it is
$(\operatorname{Im}(\bar wH))^2/|w+sH|^3$. The collinear case has
zero right side and follows by convexity.
Applying this to $w=L\pm F$ and subtracting $pt^2\|H\|^2$ supplies
a lower bound for the symmetrized actual improvement. It is integrable
because its positive integrand is bounded by a constant times $t|H|$.
This retains phase information and avoids inverse moments, but it is
not a proved positive gain for the open class.

\subsubsection{What this changes, and what it does not}
The scalar graded-noise cuts in \autoref{ch:weakk} already gave diagonal
compressions; \autoref{ch:response-stability} gave response-pair stability, and \autoref{ch:nearmax}
found its simplest tested form inactive in the old certificate.
The additions here are the complete one-fresh-coordinate quadratic block,
its off-diagonal phase information, and the quantitative identity tying
its price-level eigenspace to the actual remote cubic derivatives.
The last identity is a repackaging of stationarity and symmetry, not a
new restriction independent of both.

The proposed rigidity target is now precise: rule out an objective above
$S+99/10^7$ at $g=1/128$ for a sequence with both small response defect
and no profitable fresh-coordinate direction. There is currently no proof
of this rigidity assertion, no positive numerical measurement of the new
block on the unresolved configurations, and no improved endpoint.
Finite-dimensional nonexistence of bad stationary points would not by
itself control an increasing-dimensional sequence; the displayed error
bounds are included to make that obligation explicit.

An analogous strategy has succeeded for Gaussian noise stability of sets:
Heilman, \emph{A Variational Proof of Robust Gaussian Noise Stability},
\texttt{arXiv:2108.04950}, classifies second-variation-stable sets as
half spaces. That theorem does not apply to this complex, sign-indefinite
chaos problem. Herry--Malicet--Poly, \texttt{arXiv:2412.06749}, describes
limits of bounded-degree polynomials, but supplies neither this rigidity
theorem nor the small numerical error required here.

